\section{Related Work}

\heading{Automated Training Parallelization}
Several systems have been proposed to automate parallelization strategies for training deep learning models~\cite{alpa-osdi22,nnscaler-osdi24,unity-osdi22}. These systems perform full-scale model planning across DP, PP, and TP dimensions. Although they employ exhaustive search algorithms to generate near-optimal training configurations, the planning process often incurs substantial time overhead—requiring several minutes to complete~\cite{nnscaler-osdi24}. This makes such approaches impractical for LMM training, which demands dynamic and adaptive pipeline scheduling.

\heading{Training Systems for Multimodal Models}
Numerous system optimizations~\cite{distmm-nsdi24,optimus-atc25,disttrain-arxiv24,spindle-asplos25,orch-mllm-arxiv25,cornstarch-arxiv25,diffusionpipe-mlsys24} have been developed to address architectural and training data heterogeneity in multimodal model training. For instance, DistMM~\cite{distmm-nsdi24} tackles model and data heterogeneity in traditional CLIP-based models, where all encoders are parallel. However, this design does not generalize to LMMs, which exhibit sequential dependencies between encoders, backbones and decoders. Similarly, both Spindle~\cite{spindle-asplos25} and Optimus~\cite{optimus-atc25} target multimodal LLMs under predefined training tasks, overlooking the inherent dynamicity in LMM training pipelines.
Specifically, Spindle is designed for a static, multi-task setting where all tasks are pre-defined before training starts. This contrasts with the dynamic and flexible input-handling capabilities required by modern LMMs, which is the focus of \sys{}.
% While effective in specific contexts, these systems exhibit suboptimal performance under dynamic training scenarios.
% Nevertheless, many of their optimizations remain complementary to \sys{}'s dynamic pipeline mechanism.

\heading{Pipeline Parallelism Scheduling}
Beyond Megatron-LM's 1F1B schedule, several alternative pipeline schedules have been proposed~\cite{chimera-sc21,hanayo-sc23,zero-bubble-iclr24,graph-pipe-asplos25}. Chimera~\cite{chimera-sc21} introduces bidirectional pipelines to reduce pipeline bubbles, while zero bubble pipeline~\cite{zero-bubble-iclr24} further minimizes bubbles by decoupling backward computations into input gradient and weight gradient phases. However, these methods assume fixed pipeline stage latencies and still suffer from the dynamic imbalance problem. GraphPipe~\cite{graph-pipe-asplos25} targets models with heterogeneous modules by organizing computations into a directed acyclic graph (DAG) rather than sequential stages, enabling more flexible scheduling. \sys{}'s pipeline schedule searcher can be extended to incorporate such custom schedules like GraphPipe.

\heading{Pipeline with Variable Sequence Lengths}
Several systems address text sequence length variations in unimodal LLM training~\cite{wlb-llm-osdi25,dynapipe-eurosys24,flexpipe-atc25,byte-transformer-ipdps23,sppo-arxiv25}. For example, DynaPipe~\cite{dynapipe-eurosys24} optimizes micro-batch construction using dynamic programming for multi-task training. WLB-LLM~\cite{wlb-llm-osdi25} proposes a fine-grained per-document sharding strategy to balance workloads in context parallelism groups. FlexPipe~\cite{flexpipe-atc25} dynamically adjusts pipeline size via a live flexibility mechanism to reassign GPU workers processing short sequences to assist with longer ones. While these approaches are effective for text-based models, they do not handle multi-modal data. In contrast, our approach supports multiple modalities within a single pipeline.

% \section{Related Work}

% \heading{Automated Training Parallelization}
% Several systems enable automated parallelization for training deep learning models.
% Notable examples include Alpa~\cite{alpa-osdi22}, nnScaler~\cite{nnscaler-osdi24} and Unity~\cite{unity-osdi22}.
% These systems perform full-scale model planning across DP, PP, and TP dimensions.
% While they employ exhaustive search algorithms to generate near-optimal training configurations,
% the planning process often requires several minutes to complete~\cite{nnscaler-osdi24}.
% This substantial time overhead renders such approaches impractical for LMM training scenarios that demands dynamic pipeline schedules.

% \heading{Training Systems for Multimodal Models}
% Many system optimizations~\cite{distmm-nsdi24,optimus-atc25,disttrain-arxiv24,spindle-asplos25,orch-mllm-arxiv25,cornstarch-arxiv25,diffusionpipe-mlsys24} have been developed to address the architectural and training data heterogeneity inherent in multimodal model training.
% DistMM~\cite{distmm-nsdi24} address model and data heterogeneity in old-fashioned CLIP-based models where all encoders are parallel,
% which lacks the backbone models in LMMs that have sequential dependency on encoders and decoders.
% As a result, DistMM's approach does not generalize to LMM training.
% Both Spindle~\cite{spindle-asplos25} and Optimus~\cite{optimus-atc25} targets multimodal LLM with a predefined training task, overlooking the inherent dynamicity in LMM training.
% While effective in their respective domains,
% their approach shows suboptimal performance in dynamic LMM training.
% Nevertheless, these optimization techniques remain complementary to \sys{}'s dynamic pipeline mechanism.

% \heading{Pipeline Parallelism Scheduling}
% In addition to Megatron-LM's 1F1B, there are other pipeline schedules~\cite{graph-pipe-asplos25,chimera-sc21,hanayo-sc23,zero-bubble-iclr24}.
% Chimera~\cite{chimera-sc21} introduces bidirectional pipelines to reduce pipeline bubbles.
% Zero bubble pipeline~\cite{zero-bubble-iclr24} eliminate most pipeline bubbles by splitting backward computations into two parts of input gradient backward and weight gradient backward.
% However, these methods assume fixed pipeline stages latencies, and thereby still suffer from the dynamic imbalance problem.
% GraphPipe~\cite{graph-pipe-asplos25} targets deep neural network models with heterogeneous modules.
% It partitions computations into a DAG instead of sequential pipeline stages for more flexible pipeline schedules.
% \sys{}'s pipeline schedule searcher can be extended to support custom pipeline schedules like GraphPipe.

% \heading{Pipeline with Variable Sequence Lengths}
% Some systems targets addressing the dynamicity of variable text sequence lengths in unimodal LLMs~\cite{wlb-llm-osdi25,dynapipe-eurosys24,flexpipe-atc25,byte-transformer-ipdps23,sppo-arxiv25}.
% DynaPipe~\cite{dynapipe-eurosys24} optimizes micro-batch construction using a dynamic programming-based approach for multi-task LLM training.
% WLB-LLM~\cite{wlb-llm-osdi25} introduces a novel fine-grained per-document sharding strategy, ensuring each worker within a context parallelism group has an identical workload.
% FlexPipe~\cite{flexpipe-atc25} dynamically adjusts pipeline parallelism size by a live flexibility mechanism to utilize idle GPU workers with short sequences for long sequence processing.
% All of these works focus on unimodal LLM training, while our approach is capable of handling multiple modalities in a single pipeline.
